%% P0: The Invisible Civilization
%% Draft v0.1 — sections 1-3 (Introduction, Demographic Puzzle, Five Channels)
%% Target venue: Journal of Anthropological Archaeology (Elsevier, Q1, subscription route)
%% Target length: 8,000-10,000 words total; this file currently ~4,000 words in §1-3
%% Companion to: P1-core (sedimentation calibration, submitted separately to JASREP)
%% Created: 2026-04-20

\documentclass[12pt,a4paper]{article}
\usepackage[utf8]{inputenc}
\usepackage[T1]{fontenc}
\usepackage{graphicx}
\usepackage{amsmath}
\usepackage{natbib}
\usepackage{booktabs}
\usepackage{array}
\usepackage{hyperref}
\usepackage[left=2.5cm,right=2.5cm,top=2.5cm,bottom=2.5cm]{geometry}
\usepackage{lineno}
\usepackage{setspace}
\doublespacing
\linenumbers

\hypersetup{colorlinks=true,linkcolor=blue,citecolor=blue,urlcolor=blue}

\begin{document}

\title{The Invisible Civilisation:\\Five Independent Lines of Evidence for an Archaeologically Erased Pre-Hindu Nusantara}

\author{Mukhlis Amien\textsuperscript{1,*} \and Go Frendi Gunawan\textsuperscript{1}}

\date{}

\maketitle

\noindent\textsuperscript{1}Lab Data Sains, Universitas Bhinneka Nusantara, Malang, Indonesia\\
\textsuperscript{*}Correspondence: amien@ubhinus.ac.id

\begin{abstract}
The archaeological record of Java before the fourth century~CE is effectively empty. The conventional narrative treats this as a chronological fact: Javanese polities formed late and adopted Indic inscription with Indic civilisation. Demographic modelling, ecological carrying capacity estimates, and comparison with contemporaneous Austronesian societies suggest a different possibility: Java at 400~CE supported roughly one to two million inhabitants, comparable to Dvaravati Thailand and exceeding the early Philippines and Sriwijaya. If this population existed, its material record has been selectively erased. This paper assembles five independent evidence channels that converge on the existence of such an erased substrate civilisation: cumulative volcanic sedimentation rates that place pre-Hindu deposits at 4--10~m depth; demographic back-projection yielding a 1{,}000- to 7{,}000-fold gap between expected and observed settlements; deep comparison with the Philippine archaeological record, which preserves 275--340 pre-400~CE sites (55--65\% open-air) against zero in Java's volcanic interior; linguistic substrate reconstruction identifying Proto-Austronesian writing vocabulary and 438 substrate items in six Sulawesi languages; and reinterpretation of published paleogenomic data showing that ancient DNA has never been recovered from any open-air context in volcanic Java, a pattern consistent with the karst-versus-open-air preservation differential that also characterises Southeast Asian paleogenomic recovery more broadly. The five channels use different data, different methods, and different observer communities. They converge on a ``selective survival'' regime in which metal and cave-preserved material survives while organic settlements do not. We argue this regime is diagnosable computationally, falsifiable through specified subsurface predictions, and visible in living tradition (wayang, \textit{gunungan}, the Semar complex). The paper closes with a six-filter framework and eight pre-registered predictions whose outcomes can support, constrain, or falsify the argument.

\medskip
\noindent\textbf{Keywords:} archaeological taphonomy; multi-channel convergence; pre-Hindu Nusantara; invisible civilisations; selective survival; volcanic tropical archaeology; Indonesian prehistory
\end{abstract}


%% ================================================================
\section{Introduction: The Dwarapala's Half-Buried Body}

In 1803, the Dutch colonial administrator Nicolaus Engelhard encountered a striking scene at the ruins of Candi Singosari in East Java. Two enormous stone guardian statues, each 370~cm tall and carved from monolithic andesite in 1268~CE during the reign of King Kertanegara, stood with approximately half their bodies below the surrounding ground surface. They had not been deliberately interred. Gunung Kelud had erupted about twenty times in the intervening 535 years, and the landscape had risen around them at approximately 3.5~mm per year \citep{kinney2003}. The statues were visible only because of their monumental scale. Structures built of wood, brick, and thatch in the same period left no comparable marker above the accumulating sediment.

This observation opens a question that has not been pursued systematically in Indonesian archaeology. If a landscape buries stone monuments to half their height within five centuries, what does it do to the organic architecture that housed most of its inhabitants? The rate Engelhard's observation implies, when extended with three additional calibration sites from the Merapi volcanic system (Sambisari, Kedulan, Kimpulan; 9th century~CE), yields a cross-system mean of approximately 4.4~mm/yr \citep[hereafter P1-core]{amien_p1core}. Applied to the pre-Hindu horizon, this rate places archaeological deposits from approximately 400~CE at depths of 3.9--10.1~m in Java's volcanic basins---beyond the reach of surface survey and at or exceeding the effective limit of conventional ground-penetrating radar.

The calibration itself is geomorphological, not archaeological. It cannot demonstrate that pre-Hindu settlements exist at those depths. It can only establish that the surface record of volcanic Java has a detection horizon that cuts off sharply within the historical period. What lies below that horizon is the subject of this paper.

The conventional reading of Java's early archaeological record is developmental: Hindu-Buddhist kingdoms crystallised around the 4th century~CE, adopted Indic inscription technology, and left the epigraphic and monumental record from which subsequent polities are reconstructed \citep{coedes1968,miksic2004}. Before this threshold: silence. The silence has typically been interpreted as absence. Kutai Martadipura in non-volcanic East Kalimantan, dated to approximately 400~CE through its Yupa inscriptions \citep{vogel1918}, is treated as Indonesia's oldest known polity. Javanese state formation is dated from the 4th or 5th century~CE inscriptions of Tarumanagara. The possibility that a comparable or larger population existed on volcanic Java before these dates, but left no surface-recoverable material record, has been considered implicitly impossible---or perhaps, more often, simply not considered.

This paper argues that the possibility is both plausible and evidenced. Not through any single discovery, but through the convergence of five independent lines of evidence: sedimentation-rate calibration, demographic back-projection, comparative archaeology against the Philippine record, linguistic substrate reconstruction, and reinterpretation of published paleogenomic data. Each channel addresses a different aspect of the problem; each uses data collected by a different research community; each was not designed with VOLCARCH's thesis in mind. What they share is a single pattern: volcanic Java's pre-400~CE material record is quantitatively incompatible with what Austronesian prehistory, ecological carrying capacity, linguistic depth, and genetic continuity would independently predict.

We develop the argument in three parts. Section~2 establishes the demographic puzzle: how many people should have lived on Java at 400~CE, and what archaeological footprint should they have left? Section~3 presents the five evidence channels, each with its own method, result, independence assessment, and limitation. Section~\ref{sec:selective_survival} introduces the concept of ``selective survival'' that organises the positive evidence---bronze drums, megaliths, wayang---alongside the negative. Section~\ref{sec:layers} develops a six-filter framework linking the evidence channels into a single cascade of archaeological invisibility. Section~\ref{sec:predictions} presents eight pre-registered predictions whose outcomes can falsify, support, or constrain the framework. Section~\ref{sec:limits} catalogues limitations honestly. Section~\ref{sec:conclusion} concludes with methodological implications for Island Southeast Asian archaeology.

A note on what this paper does not claim. We do not claim to have discovered a lost civilisation. We do not claim that any specific location contains pre-Hindu remains. We do not claim that the demographic estimates are precise---they span an order of magnitude, and we present them as bounds rather than measurements. We do claim that the five channels, when read together, make a purely developmental explanation of Java's pre-400~CE record quantitatively untenable, and that the alternative explanation---archaeological invisibility produced by compound taphonomic processes---is both testable and consistent with what positive evidence does survive.


%% ================================================================
\section{The Demographic Puzzle}
\label{sec:puzzle}

\subsection{What the record shows}

Java's pre-400~CE archaeological record is minimal. No inscriptions predate the 4th-century Tarumanagara stones in non-volcanic West Java \citep{vogel1918,coedes1968}. No monumental architecture predates the 7th-century Dieng temples in the Kedu plain. Open-air habitation sites identified as pre-Hindu exist in the low single digits, most with contested or imprecise dating, and all outside Java's central volcanic interior. A compilation of 666 known archaeological sites in East Java \citep[P1-core]{amien_p1core} yields fewer than five percent attributed to any period earlier than the 10th century~CE, and zero from Java's volcanic interior that can be confidently dated before 400~CE.

This record is concentrated in three ways. First, geographically: known sites cluster within 50~km of major volcanic centres in the Singosari-Majapahit heartland around Malang, Mojokerto, and the Prambanan plain. Second, materially: the record is dominated by stone architecture (\textit{candi}) that weighs thousands of tonnes and protrudes through accumulated sediment. Third, chronologically: it favours the 8th through 15th centuries~CE, the periods when Hindu-Buddhist polities commissioned stone construction. Each concentration reflects a different bias. Together they produce a record that maps monument durability, survey effort, and genre-specific inscriptional practice---not settlement distribution.

\subsection{What Austronesian ecology predicts}

A population estimate for 400~CE Java can be derived from three independent approaches, all of which converge on the same order of magnitude.

The ecological carrying capacity approach applies ethnographically documented densities for tropical Austronesian agricultural systems---approximately 5 to 30 persons per km\textsuperscript{2} for swidden and early wet-rice cultivation \citep{bellwood2017,kirch2000,baylisssmith1980}---to Java's habitable lowlands. Using a terrain-suitability threshold that excludes slopes steeper than roughly 15 degrees and elevations above 1{,}000~m, habitable area is approximately 65{,}000~km\textsuperscript{2}, yielding a carrying-capacity range of 325{,}000 to 1{,}950{,}000 inhabitants.

Demographic back-projection from early-modern anchors provides an independent check. Reid's compilation places Java's 1600~CE population at approximately four to five million \citep{reid1988}. Plausible Malthusian growth rates (0.05--0.15\% per annum averaged across the 1200 intervening years, accounting for eruption disruptions, trade-cycle fluctuations, and the 14th-century Black Death analogues) imply a 400~CE baseline on the order of 600{,}000 to two million---consistent with the carrying-capacity result without sharing any of its assumptions.

Comparative scaling against contemporaneous Austronesian societies places Java at or above the regional median. Dvaravati Thailand supported 300{,}000--500{,}000 inhabitants during roughly the same period \citep{higham2014}; the early Philippines 200{,}000--500{,}000 \citep{junker1999}; early Sriwijaya 50{,}000--150{,}000 \citep{manguin2004}. Java, with the densest volcanic-soil fertility in the archipelago, would be an outlier if its 400~CE population fell below the Thai figure.

The three approaches converge on a central estimate of one to two million inhabitants at 400~CE, with a conservative lower bound near 600{,}000 and an upper bound near three million. This convergence is not coincidental; the approaches share no data and do not use each other's intermediate quantities. We adopt 1--2 million as a working figure, with sensitivity to the lower bound tested below.

\subsection{What the gap implies}

At standard Austronesian village sizes of 100 to 200 persons, a population of one to two million implies between 5{,}000 and 20{,}000 settlements. The observed count of pre-400~CE settlement sites in volcanic Java's interior is zero to three, depending on dating confidence. The discrepancy, computed as expected settlements divided by observed sites, spans approximately three orders of magnitude under plausible parameter choices. Under the minimal-population scenario (590{,}000 inhabitants, maximum village size of 200, and liberal counting of 3 sites) the gap is approximately 1{,}000-fold. Under the moderate scenario (1{.}9 million inhabitants, same village and site parameters) the gap is approximately 3{,}200-fold (E108). Under the maximum scenario (3{.}9 million inhabitants, same parameters) the gap reaches approximately 6{,}500-fold. Across the full range of scenario parameters tested, the gap remains above 1{,}000-fold.

Three potential explanations present themselves. First, the population estimate could be wrong. Second, the settlements could have taken forms that leave no archaeological trace regardless of taphonomy. Third, the archaeological record could be biased against detecting them.

The first explanation is difficult to sustain. To close the gap through lower population alone would require 400~CE Java to have fewer than 2{,}000 inhabitants---a figure inconsistent with every independent estimation method, every comparative analogue, and every demographic back-projection that reaches the period. Java would need to be empty, in an archipelago where neighbouring Austronesian societies with far poorer soils supported hundreds of thousands.

The second explanation---that settlements left no archaeological trace---has surface plausibility. Organic construction using bamboo, wood, palm leaf, and thatch is undetectable after a few centuries in the tropical climate even without volcanic burial \citep{erasmus2005}. However, the organic-architecture explanation alone predicts that non-volcanic tropical Austronesian settings should also show similar archaeological absence. They do not. The Philippine record, discussed in Section~\ref{sec:channels}, preserves hundreds of open-air pre-400~CE sites across varied environments. The absence, in other words, is specifically concentrated in Java's volcanic terrain, not distributed across all comparable tropical Austronesian landscapes.

The third explanation---taphonomic bias in the archaeological record---is the hypothesis this paper develops. Not as a replacement for the first two (population was certainly smaller than the 20{,}000-settlement upper bound, and organic architecture certainly contributes to invisibility) but as the load-bearing factor that, together with the others, makes the discrepancy coherent. The remainder of this paper assembles the five evidence channels that support this interpretation.


%% ================================================================
\section{Five Independent Evidence Channels}
\label{sec:channels}

The argument proceeds through five channels, each drawing on different data, different methods, and different observer communities. We assess each separately, then discuss their combined implications. The channels are: cumulative volcanic sedimentation (\S\ref{sec:ch1}); comparative archaeology against the Philippine record (\S\ref{sec:ch2}); linguistic substrate reconstruction (\S\ref{sec:ch3}); paleogenomic reinterpretation (\S\ref{sec:ch4}); and colonial archive text mining (\S\ref{sec:ch5}).

Table~\ref{tab:channels} summarises the channels, their primary evidence, and their independence from one another.

\begin{table}[t]
\caption{The five evidence channels, primary sources, and independence assessment. ``Independent of'' indicates channels that share no primary data or method. Note: all channels that invoke population or comparative estimates draw ultimately on the same ethnographic and historical-demographic literature base (principally \citealp{bellwood2017,reid1988,kirch2000,higham2014}); the channels are methodologically and evidentially diverse rather than fully epistemically independent.}
\label{tab:channels}
\begin{tabular}{p{2.5cm}p{4.5cm}p{4cm}p{3cm}}
\toprule
Channel & Primary data & Observer community & Independent of \\
\midrule
Sedimentation (\S\ref{sec:ch1}) & 4 calibration sites + 51 eruption--site pairs & Geoarchaeologists, volcanologists & 2, 3, 4, 5 \\
Philippines comparison (\S\ref{sec:ch2}) & ${\sim}275{-}340$ Philippine pre-400 CE sites & Southeast Asian archaeologists & 1, 3, 4 (partial 5) \\
Linguistic (\S\ref{sec:ch3}) & ABVD wordlists, OJW Wordnet, 438 substrate items & Historical linguists, computational linguists & 1, 2, 4 (partial 5) \\
Genomic (\S\ref{sec:ch4}) & Published aDNA + WGS datasets & Paleogeneticists & 1, 2, 3, 5 \\
Colonial archive (\S\ref{sec:ch5}) & 16 OV volumes, 1{,}768 Delpher articles & Colonial-era Dutch and Indonesian fieldworkers & 3, 4 (partial 1, 2) \\
\bottomrule
\end{tabular}
\end{table}

\subsection{Volcanic sedimentation rates and the detection horizon}
\label{sec:ch1}

The first channel is the one whose methodological detail is developed in full in a companion calibration paper \citep{amien_p1core}. Here we summarise its result for the synthesis.

Four archaeological sites with independently documented construction dates and measured burial depths were compiled from two volcanic systems. The Dwarapala statues of Singosari (1268~CE, Kelud system, 185~cm of accumulation over 535 years) yield a rate of 3.5~mm/yr, internally consistent with Kelud's approximately twenty eruptions in that period. Three Merapi-system temples---Sambisari (835~CE, 500--650~cm), Kedulan (869~CE, 600--700~cm), Kimpulan (900~CE, 270--500~cm)---yield rates of 4.4--5.7, 5.3--6.2, and 2.4--4.5~mm/yr respectively. The four-site mean is approximately 4.4~mm/yr across 350~km of Java and two geologically distinct volcanic basins.

Two independent validation datasets test this rate. First, 51 eruption--archaeological site pairs, compiled from colonial and volcanological literature spanning 12 volcanic systems, yield a mean measured burial depth of 3.41~m. Applying the four-site calibration rate to a 760-year horizon predicts 3.3~m---an agreement within 4\% that was not a design target of either compilation. Second, a larger compilation of 363 site-depth observations from East Java archaeological service reports (P1-core, E075) yields a Pearson correlation of $r=0.951$ between predicted and observed burial depths.

Applied as a first-order linear projection, the calibration rates place pre-400~CE deposits in the Malang basin at 3.9--10.1~m depth, Kanjuruhan-era (760~CE) deposits at 3.0--7.9~m, and Singosari-era (1268~CE) deposits at 1.8--4.7~m. The effective range of ground-penetrating radar in volcanic sediment is 2--5~m \citep{conyers2004,goodman2013}; the effective range of systematic surface survey is below 1~m. The pre-400~CE horizon is therefore beneath both techniques' practical reach, except through construction-triggered exposure or deep coring.

The channel establishes a detection horizon, not a buried record. It cannot show that settlements exist at the projected depths; it can show that the surface archaeological record of Java's volcanic interior cuts off within the historical period. Whether the cut-off reflects taphonomic invisibility or cultural absence is a question the other four channels address.

This channel is falsifiable. Systematic deep coring or construction-triggered exposure to depths of 4--10~m in volcanic-interior target zones that yields zero anthropogenic material, across a statistically adequate sample, would reject the detection-horizon interpretation. Rates systematically lower than 4.4~mm/yr in additional calibration sites would reduce the projected pre-400~CE depth and invalidate the claim that this horizon exceeds conventional survey reach.

\subsection{The Philippines natural experiment}
\label{sec:ch2}

The Philippines shares Java's major cultural and environmental parameters: Austronesian heritage, tropical climate, island-archipelagic geography, active volcanism, and organic building tradition. It differs in three relevant ways: lower volcanic density (approximately 0.07 active volcanoes per 1{,}000~km\textsuperscript{2} versus Java's 0.35), abundant karst providing natural preservation contexts, and smaller pre-colonial population by all estimates. If volcanic burial were not a significant archaeological factor, the Philippine pre-400~CE record should be proportionally thinner than Java's given its smaller population; if it is, the record should be proportionally richer given its preservation advantages.

A deep composition analysis of the Philippine archaeological record \citep[E201]{amien_e201} identifies approximately 275--340 documented pre-400~CE sites across site types: around 30--50 cave habitation sites (Tabon, Callao, Balobok, Pilanduk, Bilat), approximately 170 open-air Paleolithic sites (Cagayan Valley alone documents 168+), 30+ Neolithic shell middens (the Lal-lo/Gattaran complex), 20--30 open Neolithic settlement sites (Andarayan, Batanes jade sites), 15--20 Metal Age cave burials, 5--10 Metal Age open burials, and 5--10 rock art sites. The site-type breakdown is 55--65\% open-air, 35--45\% cave-based---a diverse record, not a cave-dominated one.

Material culture representation follows the same pattern. The Philippine pre-400~CE record includes abundant pottery (from approximately 4{,}000~BP onward), stone tool assemblages (700{,}000~BP through Neolithic), metal objects (500~BCE onward, including gold, bronze, and iron), tens of thousands of jade artifacts from the Batangas culture (2000--1500~BCE), hundreds of human burials in both cave and open contexts, evidence of rice from 1500--1400~BCE (Andarayan), and domesticated pig remains from 2500--2200~BCE (Nagsabaran). Volcanic Java's pre-400~CE record contains none of these categories.

Critically, Philippine volcanic zones themselves preserve archaeological material. Batangas---adjacent to Taal volcano, which has erupted 39 or more times historically---preserves the Calatagan burial grounds, jade culture sites, and coastal shell middens. Batanes sites buried under Iraya ash were recovered by systematic excavation \citep{bellwood_dizon2013}. The site-forming potential of volcanic burial, demonstrated in the Philippines, is recovery-dependent rather than preservation-destroying.

Volcanic Java's zero pre-400~CE open-air sites stand against documented pre-400~CE sites in Philippine volcanic zones. Within Java itself, non-volcanic West Java preserves pre-Hindu material: the Buni culture (c.~400~BCE to 500~CE) on the northern coast and Batujaya's 1st--3rd~century~CE prehistoric burial layer beneath the later Tarumanagara temple complex demonstrate that open-air pre-400~CE archaeology is recoverable where volcanic sedimentation is absent \citep{manguin2011,wibisono1994}. The within-island contrast---non-volcanic Java preserves, volcanic Java does not---constrains the possible explanations. The remaining factors are (a) Philippines' more abundant karst, which provides natural preservation capsules bypassing other taphonomic filters; (b) Java's more continuous volcanic chain, producing more extensive tephra blankets; and (c) the assumption in Indonesian archaeology that the volcanic interior is empty, which has directed survey effort elsewhere and become self-fulfilling.

The Philippines comparison provides the strongest single external constraint on the Java thesis. It controls for climate, Austronesian heritage, tropical ecology, organic architecture tradition, and tropical organic decay. What remains as differential is volcanism, karst availability, and survey attention. The zero-versus-hundreds site-count contrast is not reducible to any single one of these factors alone.

This channel is falsifiable. Discovery of open-air pre-400~CE sites in Java's volcanic interior, through systematic survey targeting predicted locations, would weaken or eliminate the within-island contrast. Recovery of organized pre-Hindu Philippine sites at sedimentation rates comparable to Java's 4.4~mm/yr would reduce the cross-regional control.


%% ================================================================
% Remaining channels (§3.3-3.5) + §4-9 continued in separate file/section.
% This file ends here as v0.1 deliverable for Pak Amien review.
% Approximate word count of §1-3 so far: ~3,500 words.

\vspace{1em}
\noindent\textit{[End of v0.1 draft. Remaining channels (linguistic, genomic, colonial archive), selective survival section, six-layer framework, predictions, limitations, and conclusions to be drafted in v0.2.]}

\bibliographystyle{apalike}
\bibliography{references}

\end{document}
